\section{Discussion}
\label{sec:discussion}

This section addresses limitations and design boundaries. We follow
the three-part structure used in companion work
\cite{blackrim_moa_paper}: where the design pays a cost, why we paid
it anyway, and when the design will not work.

\subsection{The cost of the false-negative-only contract}
\label{sec:discussion:falsenegative}

\Cref{def:false_negative} forbids the compile gate from accepting
diffs that quietly break the program but explicitly permits it to
reject diffs that would in fact be safe. The cost is friction: an
LLM agent proposing a structural change that the gate rejects has to
re-plan, with no guarantee that its next attempt will fare better.
In a session of $K$ refactor attempts, the expected number of false
negatives at rate $\phi$ is $K \phi$; at $K = 20$ and $\phi = 0.05$
that is one re-plan per session, which is acceptable.

The trade-off was deliberate. The opposite polarity --- a gate that
admits anything compile-clean and rejects only what the native
syntax checker catches --- would let through a class of subtle
refactor bugs we cannot tolerate: imports removed because the only
reference was inside a renamed identifier, function signatures
updated at the declaration but not at all call sites, type aliases
quietly broken. The native syntax check is necessary but not
sufficient; the LSP-diagnose step (\cref{alg:gate}) adds the
sufficiency where available.

\subsection{Tree-sitter version brittleness}
\label{sec:discussion:treesitter}

\Cref{assn:significance} assumes per-language significance rules are
stable across files for a fixed grammar version. In practice,
tree-sitter grammar versions rename node-type names roughly every
9--18 months. Our per-language significance tables (\cref{app:significance})
list explicit node names; a grammar bump that renames
\texttt{function\_declaration} to \texttt{function\_def} will silently
break the rule for that construct.

The mitigations are: (a) grammar versions are pinned in
\texttt{go.mod} and the in-tree parser-registration tables;
(b) version bumps require updating the per-language node-kind tables
in lockstep with re-running the acceptance fixtures, which encode
canonical Markdown outputs for each language;
(c) the universal fallback (line-span $\geq 10$ or named construct)
provides a partial floor when a recognised construct is silently
missed.

The brittleness is real and not fully addressed by the design. A
fully grammar-independent solution would require a query layer over
tree-sitter that abstracts node names into language-neutral
categories; we treat that as future work
(\cref{sec:future:tier2}). The empirical question is whether the
grammar-bump cadence is faster than our acceptance-fixture
maintenance cadence; on the evidence of the past two grammar
revisions (\tothink{which revisions, when, what broke}), it has not
been.

\subsection{Sanitisation as fidelity loss}
\label{sec:discussion:sanitization}

The sanitisation pass (\cref{sec:algorithm:sanitization}) is a hedge
against prompt injection from verbatim file content. It is also a
fidelity-loss surface: a legitimate package doc that happens to
contain the literal string \texttt{SYSTEM:} (e.g., a Python module
documenting an HTTP \texttt{System:} header convention) will be
flagged with a \texttt{[sanitised]} prefix. The flag is visible to
the agent and to a human reviewer; nothing is dropped, only marked.

The conservative-biased choice was deliberate. The alternative ---
trust verbatim content from any source --- would expose every agent
that auto-prepends an outline to a class of supply-chain attacks: a
malicious vendored package whose top-of-file comment is a
prompt-injection payload would be ingested into every reading
agent's context. We pay the legitimate-prose-misflagged cost to rule
out that class entirely.

\subsection{When this won't work}
\label{sec:discussion:limits}

The design assumes textual source code whose structure is meaningfully
captured by an AST. It will not work for:

\begin{itemize}[leftmargin=*]
  \item \textbf{Non-textual code} (Jupyter notebooks, binary
        serialised models, JSON-defined low-code workflows). The
        AST view is null; agents must fall back to format-specific
        tooling.
  \item \textbf{Generated code} (protoc output, ORM-generated
        models, transpiled bundles). The outline is technically
        valid but uninformative; the \texttt{outline:skip} escape
        hatch is the operator's lever.
  \item \textbf{Constant-table--dominated files} (10000-line
        embedded lookup tables, generated language-data files). The
        AST view degenerates to one large declaration; the
        compression ratio approaches 1 and the LoC threshold
        triggers no useful outline. Operator's lever:
        \texttt{outline:skip} or per-project
        \texttt{threshold\_loc} bump.
  \item \textbf{Adversarial inputs that defeat sanitisation.}
        \Cref{sec:algorithm:sanitization} catches known
        prompt-injection patterns but is not robust against a
        determined attacker who crafts content that semantically
        induces an LLM to ignore prior instructions without using
        any of the flagged literal strings. The conservative-biased
        sanitisation is necessary but not sufficient.
\end{itemize}

\subsection{Off-policy gap and agents already in the wild}
\label{sec:discussion:offpolicy}

The end-to-end token reduction claim (\cref{sec:eval:e2e}) is
fundamentally an on-policy measurement: the agent's behaviour
changes in response to the discipline subsystem (it learns to
expect outlines, to navigate from outlines, to fetch bodies
selectively). A counterfactual estimate from logged traces
\emph{before} the discipline rollout would systematically
under-estimate the reduction because pre-rollout agents do not know
how to consume outlines efficiently.

This is structurally identical to the off-policy evaluation gap in
the companion model-advisor work \cite{blackrim_moa_paper}. The
mitigation there --- doubly-robust off-policy evaluation
\cite{dudik2014doubly} keyed by per-cell logged propensities --- is
\emph{not} directly applicable here because the policy change is in
the agent's reading strategy, not the dispatcher's routing decision.
Our analogue is paired-trace measurement on tasks run twice (once
with discipline, once without); accumulating that dataset is the
work item gating OQ-AST-5.

\subsection{Provisional inclusion of the pattern-DSL subsystem}
\label{sec:discussion:patterndsl}

The pattern-DSL authoring-velocity subsystem
(\cref{sec:system:patterndsl}) is included \emph{provisionally}: it
is the only subsystem of the five whose central empirical contribution
(intents authored per engineer-hour, with measured fidelity loss vs.
native parser walking) cannot yet be measured because the
implementation is not yet shipped. We discuss it for completeness
and because the architectural design --- a YAML grammar layered on a
choice of two execution backends (ast-grep binary or in-tree Go
fallback) --- is settled and consequential for the design's overall
shape. The paper's empirical contribution does not depend on it; if
the subsystem is delayed past the paper's revision window, the
relevant claims downgrade from ``provisional'' to ``planned'' without
disturbing the rest of the evaluation.
